\section{遍历理论在数论中的应用}

本节的目的在于证明以下的命题：
设$\{F_n\}_{n\in\mb{N}}$是斐波那契数列，$F_{n+2}=F_{n+1}+F_n,F_0=0,F_1=1$，则对于几乎所有的实数$\alpha$，$\{F_n\cdot\alpha\mod 1\}_{n\in\mb{N}}$在$[0,1)$中均匀分布，这等价于$\forall\, f\in C(S^1)(S^1\cong [0,1),S^1=\,\text{圆周}\,)$，$\frac{1}{N}\sum_{n=0}^{N-1}{f\left(F_n\cdot\alpha\right)}\to\int_{S^1}{f}\d x$.\par
其中要使用的最重要的是下列的定理：
\begin{theorem}{紧致度量空间上连续函数空间的可分性}
    设$X$是一个紧致度量空间，$C(X)=\{X\,\text{上的连续函数}\}$，定义$C(X)$上的范数为$\|f\|=\sup_{x\in X}|f(x)|$，定义$d(f,g)=\|f-g\|$，则$(C(X),\|\cdot\|)$是一个可分的度量空间，即存在$\{f_k\}_{k\in\mb{N}}\subseteq C(X)$，使得$\{f_k\}_{k\in\mb{N}}$在$C(X)$中稠密.
\end{theorem}
考虑矩阵$A=\begin{pmatrix}
    0 & 1 \\
    1 & 1
\end{pmatrix}$和$T_A:X_2\to X_2$，则$|\lambda I-A|=\begin{vmatrix}
    \lambda & -1 \\
    -1 & \lambda-1
\end{vmatrix}=\lambda(\lambda-1)=\lambda^2-\lambda-1=0$，$A$的特征值均不是单位根，因此$T_A$是遍历映射.\par
由Birkhoff遍历定理，$\forall\, f\in L^1(X_2,m_\mathrm{Leb})$，$\frac{1}{N}\sum_{n=0}^{N-1}{f \left(T^n_A(x)\right)}\to \int_{X_2}{f}\d m_\mathrm{Leb}$，$\aew x+\mb{Z}^2\in X_2\xlongrightarrow{\text{提升到}}\aew x\in\mb{R}^2$.\par

根据上面的定理，$C(X_2)\subseteq L^1(X_2,m_\mathrm{Leb})$是可分空间，设$\{f_k\}_{k\in\mb{N}}\subseteq C(X_2)$是其中的一个可数稠密子集.\par
$\forall\, k\in\mb{N}$，存在$\mb{R}^2$中的全测集$Q_k$，使得$\forall\, x\in Q_k,\frac{1}{N}\sum_{n=0}^{N-1}{f_k \left(T^n_A(x)\right)}\to \int_X{f_k}\d m_\mathrm{Leb}$，取$Q=\bigcap_{k=1}^\infty{Q_k}$，则$Q$也是$\mb{R}^2$中的全测集.\par
$\forall\, f\in C(X_2),\forall\, \varepsilon>0$，由$\{f_k\}_{k\in\mb{N}}$稠密知$\exists\, k\in\mb{N},\st d(f_k,f)<\varepsilon \Leftrightarrow \sup_{x\in X_2}{|f(x)-f_k(x)|}<\varepsilon$.\par
$\forall\, x\in Q$，$\frac{1}{N}\sum_{n=0}^{N-1}{f_k \left(T^n_A(x)\right)}\to \int_{X_2}{f_k}\d m_\mathrm{Leb}$$\Rightarrow \varlimsup_{n\to\infty}{\left|\frac{1}{N}\sum_{n=0}^{N-1}{f \left(T^n_A(x)\right)}-\int_{X_2}{f}\d m_\mathrm{Leb}\right|}\leq 2\varepsilon $.\par
令$\varepsilon\to 0$，$\frac{1}{N}\sum_{n=0}^{N-1}{f \left(T^n_A(x)\right)}\to \int_{X_2}{f}\d m_\mathrm{Leb}$.\par
接下来研究$T^n_A(x)$.按照定义$T^{n+1}_A(x)=A^{n+1}x+\mb{Z}^2$.\par
其中$A^{n+1}=\underbrace{\begin{pmatrix}
    0 & 1\\
    1 & 1
\end{pmatrix}\begin{pmatrix}
    0 & 1\\
    1 & 1
\end{pmatrix}\cdots\begin{pmatrix}
    0 & 1\\
    1 & 1
\end{pmatrix}}_{n\,\text{个}}=\begin{pmatrix}
    q_n & q_{n-1} \\
    p_n & p_{n-1}
\end{pmatrix}$，其中$q_{n+1}=a_{n+1}q_n+q_{n-1}=q_n+q_{n-1},p_{n+1}=a_{n+1}p_n+p_{n-1}=p_n+p_{n-1}$满足斐波那契数列的递推公式.\par
当$n=0$时，$\begin{pmatrix}
    0 & 1 \\
    1 & 1 \\
\end{pmatrix}=\begin{pmatrix}
    q_0 & q_{-1} \\
    p_0 & p_{-1}
\end{pmatrix}\Rightarrow q_0=0,q_{-1}=1,p_0=1,p_{-1}=1$.\par
因此对于$\{p_n\}_{n\in\mb{N}}$，有$p_n=F_{n+2}$.\par
$\forall\, x\in Q,x=\begin{pmatrix}
    \alpha \\
    \beta
\end{pmatrix}\in\mb{R}^2$，$A^{n+1}x=\begin{pmatrix}
    q_n & q_{n-1} \\
    p_n & p_{n-1}
\end{pmatrix}\begin{pmatrix}
    \alpha \\
    \beta
\end{pmatrix}=\begin{pmatrix}
    \alpha q_n+\beta q_{n-1} \\
    \alpha p_n+\beta p_{n-1}
\end{pmatrix}$.\par
取$g,h\in C(S^1)$，则$f(u,v)=g(u)\cdot h(v)\in C(X_2)$.\par
则$\forall\, x\in\begin{pmatrix}
    \alpha \\
    \beta
\end{pmatrix}\in Q,\frac{1}{N}\sum_{n=0}^{N-1}{f g(\alpha q_n+\beta q_{n-1})h(\alpha p_n+\beta p_{n-1})}\to (\int_{S_1}{g}\d x)(\int_{S^1}{h}\d x)$.\par
令$g=\bf{1}$，则$\forall\, h\in C(S^1),\forall\, x=\begin{pmatrix}
    \alpha \\
    \beta
\end{pmatrix}\in Q,\frac{1}{N}\sum_{n=0}^{N-1}{h(\alpha p_n+\beta p_{n-1})}\to \int_{S^1}{h}\d x$.\par
$\Rightarrow \forall\,\begin{pmatrix}
    \alpha \\
    \beta
\end{pmatrix}\in Q,\{\alpha p_n+\beta p_{n-1}\}_{n\in\mb{N}}$在$[0,1)$中均匀分布.\par
由于$p_n=F_{n+2}$，所以$\{\alpha F_{n+1}+\beta F_n\}_{n\in\mb{N}}$也在$[0,1)$中均匀分布.\par
我们知道$F_n$的通项公式为$F_n=c(\lambda_1^n+\lambda_2^n)$，其中$|\lambda_1|>1,|\lambda_2|<1$.\par
那么$\setjot[-4]\begin{aligned}[t]
    \alpha F_{n+1}+\beta F_n &=\alpha c(\lambda_1^{n+1}+\lambda_1^{n+1})+\beta c(\lambda_1^n+\lambda_2^n)=\lambda_1^n(c\alpha\lambda_1+c\beta)+(\lambda_2^n\cdot c\beta+c\alpha\cdot \lambda_2^{n+1}) \\
    &= F_n(\alpha \lambda_1+\beta)+\tilde{c}\lambda_2^n(\tilde{c}\,\text{为某个常数}) \\
    &=\tilde{\alpha}F_n+\tilde{c}\lambda_2^n
\end{aligned}$\par
因此$\frac{1}{N}\sum_{n=0}^{N-1}{h(\tilde{\alpha}F_n)}=\frac{1}{N}\sum_{n=0}^{N-1}{\left(h(\alpha F_{n+1}+\beta F_n)-\tilde{c}\lambda_2^n\right)}=\frac{1}{N}\sum_{n=0}^{N-1}{\left(h(\alpha F_{n+1}+\beta F_n)\right)}-\frac{\tilde{c}}{N}\sum_{n=0}^{N-1}{\lambda_2^n}$.\par
当$N\to\infty$时，$\frac{1}{N}\sum_{n=0}^{N-1}{\left(h(\alpha F_{n+1}+\beta F_n)\right)}\to\int_{S_1}{h}\d x$，$\frac{\tilde{c}}{N}\sum_{n=0}^{N-1}{\lambda_2^n}=\tilde{c}\cdot\frac{1-\lambda_2^N}{N(1-\lambda_2)}\stackrel{|\lambda_2|<1}{\longrightarrow} 0$.\par
故$\forall\, h\in C(S^1),\frac{1}{N}\sum_{n=0}^{N-1}{h(\tilde{\alpha}F_n)}\to \int_{S^1}{h}\d x\Rightarrow \{\tilde{\alpha}F_n\}_{n\in\mb{N}}$在$[0,1)$中均匀分布.\par
最后只需证明上述$\tilde{\alpha}$的全体构成一个全测集，即$m_\mathrm{Leb}\left(\left\{\tilde{\alpha}\left|\tilde{\alpha}=\alpha\lambda_1+\beta,\begin{pmatrix}
    \alpha \\
    \beta
\end{pmatrix}\in Q\right.\right\}\right)=1$.

\begin{proof}
$Q\subseteq \mb{R}^2$为全测集，由Fubini定理，存在$A\subseteq\mb{R}$的全测集，使得$\forall\,\alpha\in A$，其“垂直切片”$Q_\alpha:=\{\beta\in\mb{R}:(\alpha,\beta)\in Q\}$
在$\mb{R}$中也是全测集.\par
固定任意$\alpha_0\in A$，令$E:=\{\tilde{\alpha}\in\mb{R}\mid \tilde{\alpha}=\alpha_0\lambda_1+\beta,\ \beta\in Q_{\alpha_0}\}=\alpha_0\lambda_1+Q_{\alpha_0}$. \par
由Lebesgue测度的平移不变性，$E$在$\mb{R}$中是全测集.\par
又$E\subseteq\{\tilde{\alpha}\in\mb{R}\mid \tilde{\alpha}=\alpha\lambda_1+\beta,(\alpha,\beta)\in Q\}$，
故$\{\tilde{\alpha}\mid \tilde{\alpha}=\alpha\lambda_1+\beta,(\alpha,\beta)\in Q\}$也是$\mb{R}$中全测集.\par
\end{proof}
至此，我们就证明了本节一开始提出的命题：
\begin{proposition}
    设$\{F_n\}_{n\in\mb{N}}$为斐波那契数列，则对于几乎所有的实数$\alpha$，$\{F_n\cdot\alpha\mod 1\}_{n\in\mb{N}}$在$[0,1)$中均匀分布；等价地，$\forall\, f\in C(S^1)$，$\frac{1}{N}\sum_{n=0}^{N-1}{f\left(F_n\cdot\alpha\right)}\to\int_{S^1}{f}\d x$.
\end{proposition}

本节的第二个例子是下面这个定理：
\begin{theorem}
    设$\a\in\mb{R}\setminus\mb{Q}$，$p(x)=\alpha x^n+a_{n-1}x^{n-1}+\cdots+a_1x+a_0(a_i\in\mb{R})$，则$p(x)\mod 1$在$[0,1)$中均匀分布.
\end{theorem}
根据均匀分布的等价条件，我们实际上只需要证明$\forall\, k\in\mb{Z}\setminus\{0\},\frac{1}{N}\sum_{n=0}^{N-1}{\cos{2k\pi p(n)}}\to 0$且$\frac{1}{N}\sum_{n=0}^{N-1}{\sin{2k\pi p(n)}}\to 0$.\par
这等价于证明$\frac{1}{N}\sum_{n=0}^{N-1}{\eu^{2k\pi \iu p(n)}}\to 0(\forall\, k\in\mb{Z}\setminus\{0\})$.\par
对于上面的定理，如果$p(n)=\alpha n+a_0$，它在$[0,1)$上均匀分布，因而$\frac{1}{N}\sum_{n=0}^{N-1}{\eu^{2k\pi \iu(\alpha n+a_0)}}\to 0$.为了将其余次数的多项式化规到这一情形，我们需要下述引理.\par
\begin{lemma}
    设$\{a_n\}_{n\in\mb{N}}$是$\mb{C}$中的有界序列，且满足$\forall\, h\in\mb{N}_{>0}$，$\frac{1}{N}\sum_{n=0}^{N-1}{\overline{a_{n+h}}a_n}\to 0$，那么$\frac{1}{N}\sum_{n=0}^{N-1}{a_n}\to 0$.
\end{lemma}
\begin{proof}
    记$b_n=\frac{1}{H}\sum_{h=0}^{H-1}{a_{n+h}}(\forall\, H\in\mb{N})$，则$b_n$实际上就是从$a_n$开始的长度为$H$的$a_i$序列的平均值.\par
    $\setjot[-2]\begin{aligned}[t]
        \left|\frac{1}{N}\sum_{n=0}^{N-1}a_n-\frac{1}{N}\sum_{n=0}^{H-1}b_n\right|&\leq \frac{1}{N}\left(\sum_{n=0}^{N-1}|a_1|+\sum_{n=N-H}^{N-1}|a_n|\right)+\frac{1}{N}\left(\sum_{n=0}^{H-1}|a_1|+\sum_{n=N}^{N+H-1}|a_n|\right) \\
        &\leq \frac{4H}{N}M\;(|a_i|\leq M)
    \end{aligned}$\par
    对于固定的$H\in\mb{N}$，$\lim_{n\to\infty}\left|\frac{1}{N}\sum_{n=0}^{H-1}a_n-\frac{1}{N}\sum_{n=0}^{N-1}{b_n}\right|=0$.\par
    由均值不等式可知$\left|\frac{1}{N}\sum_{n=0}^{N-1}b_n\right|\leq \sqrt{\frac{1}{N}\sum_{n=0}^{N-1}|b_n|^2}$.\par
    而$\setjot[-2]\begin{aligned}[t]
        \frac{1}{N}\sum_{n=0}^{N-1}|b_n|^2 &=\frac{1}{N}\sum_{n=0}^{N-1}\left|\frac{1}{N}\sum_{h=0}^{H-1}a_{n+h}\right|^2=\frac{1}{N}\frac{1}{H^2}\sum_{n=0}^{N-1}\left(\sum_{i=0}^{H-1}\overline{a_{n+i}}\right)\left(\sum_{j=0}^{H-1}{a_{n+j}}\right) \\
        &=\frac{1}{N}\frac{1}{H^2}\sum_{n=0}^{N-1}\sum_{i=0}^{H-1}\sum_{j=0}^{H-1}\overline{a_{n+i}}a_{n+j}=\frac{1}{H^2}\sum_{0\leq i,j\leq H-1}\left(\frac{1}{N}\sum_{n=0}^{N-1}{\overline{a_{n+1}}a_{n+j}}\right) \\
        &=\frac{1}{H^2}\sum_{0\leq i=j\leq H-1}\left(\frac{1}{N}\sum_{n=0}^{N-1}{\overline{a_{n+1}}a_{n+j}}\right)+\frac{1}{H^2}\sum_{\substack{0\leq i,j\leq H-1\\ i\neq j}}\left(\frac{1}{N}\sum_{n=0}^{N-1}{\overline{a_{n+1}}a_{n+j}}\right)
    \end{aligned}$\par
    固定$H\in\mb{N}$，则$\varlimsup_{N\to\infty}{\frac{1}{N}\sum_{n=0}^{N-1}{|b_n|^2}}\leq\frac{M}{H}\quad \left(|a_i|\leq M,\forall\, i\in\mb{N}\right)$.\par
    $\Rightarrow \varlimsup_{N\to\infty}\left|\frac{1}{N}\sum_{n=0}^{N-1}{b_n}\right|\geq\frac{M}{\sqrt{H}}$.\par
    令$H\to\infty$，得$\setjot[-2]\begin{aligned}[t]
        \varlimsup_{N\to\infty}\left|\frac{1}{N}\sum_{n=0}^{N-1}a_n\right|&\leq\varlimsup_{N\to\infty}\left|\frac{1}{N}\sum_{n=0}^{N-1}{a_n}-\frac{1}{N}\sum_{n=0}^{N-1}b_n\right|+\varlimsup_{N\to\infty}\left|\frac{1}{N}\sum_{n=0}^{N-1}b_n\right| \\
        &\leq 0+\frac{M}{\sqrt{H}} \to 0\;(H\to\infty)
    \end{aligned}$\par
\end{proof}

借助这个引理，我们就可以证明上述的定理，先以三次多项式为例.
\begin{instance}
    若$p(n)=\alpha n^3+n$，需要证明$\frac{1}{N}\sum_{n=0}^{N-1}{\eu^{2k\pi \iu p(n)}}\to 0\;(\forall\, k\in\mb{Z}\setminus\{0\})$.\par
    设$a_n=\eu^{2k\pi \iu p(n)}$，即证$\frac{1}{N}\sum_{n=0}^{N-1}{\overline{a_{n+h}}a_n}\to 0\;(\forall\, h\in\mb{N}_{>0})$. \par
    其中$\overline{a_{n+h}}a_n=\eu^{-2k\pi \iu p(n+h)}\eu^{2k\pi \iu p(n)}=\eu^{-2k\pi \iu \left(\alpha(3n^2h+3nh^2+h^3)+h\right)}$.\par
    记$b_n=\eu^{-2k\pi \iu \left(\alpha(3n^2h+3nh^2+h^3)+h\right)}$，需要证明$\frac{1}{N}\sum_{n=0}^{N-1}{b_n}\to 0\;(\forall\, h\in\mb{N}_{>0})$.\par
    $\overline{b_{n+s}}b_n=\eu^{-2k\pi\iu \left(\alpha(3h(2ns+s^2)+3sh^2)\right)}$，这是关于$n$的一次多项式，因此$\frac{1}{N}\sum_{n=0}^{N-1}{\overline{b_{n+s}}b_n}\to 0\;(\forall\, s\in\mb{N}_{>0})$.\par
    因此$\frac{1}{N}\sum_{n=0}^{N-1}b_n=\frac{1}{N}\sum_{n=0}^{N-1}{\overline{a_{n+h}a_n}}\to 0\;(\forall\,h\in\mb{N}_{>0})$. \par
    $\Rightarrow \frac{1}{N}\sum_{n=0}^{N-1}{a_n}\to 0\;(\forall\, k\in\mb{Z}\setminus\{0\})$.
\end{instance}
类似上面的例子，对于一般的多项式$p(n)$，只需要对$\mathrm{deg}\, p$用数学归纳法即可.另外，上述定理实际上可以加强为：
\begin{theorem}{}
    $p(x)=a_nx^n+a_{n-1}x^{n-1}+\cdots+a_1x+a_0$，若$a_i(1\leq i\leq n)$不全为有理数，则$p(x)\mod 1$在$[0,1)$中均匀分布.
\end{theorem}
\begin{proof}
    我们对多项式的次数$n$使用第二数学归纳法.\par
    当$n=1$时，$p(x)=a_1x+a_0$，其中$a_1$是无理数，成立.\par
    假设定理对所有次数小于$n$的多项式成立，考虑$n$次多项式$p(x)=a_nx^n+\cdots+a_0$.\par
    设$m$是使得$a_m$为无理数的最大下标($1\leq m\leq n$).\par
    要证$\frac{1}{J}\sum_{j=0}^{J-1}{\eu^{2k\pi \iu p(j)}}\to 0\;(\forall\, k\in\mb{Z}\setminus\{0\})$，设$b_j=\eu^{2k\pi \iu p(j)}$，即证$\frac{1}{J}\sum_{j=0}^{J-1}{\overline{b_{j+h}}b_j}\to 0\;(\forall\, h\in\mb{N}_{>0})$，其中$\overline{b_{j+h}}b_j=\eu^{-2k\pi\iu(p(j+h)-p(j))}$. \par
    记$q_h(j)=p(j+h)-p(j)$，$q_h(j)=\sum_{k=1}^n a_k ((j+h)^k - j^k)$.\par
    考虑$q_h(j)$中$j^{m-1}$的系数，它来自于$q_h(x)$中$k\geq m$的项展开式中.\par
    由二项式定理$a_k((j+h)^k - j^k) = a_k\left(kj^{k-1}h + \binom{k}{2}j^{k-2}h^2 + \cdots\right)$.\par
    因此，$q_h(j)$中$j^{m-1}$的系数为$c_{m-1} = mha_m + \sum_{k=m+1}^n a_k \cdot \binom{k}{k-m+1}h^{k-m+1}$.\par
    由$\forall\,k>m$，$a_k\in\mb{Q}$且$h\in\mb{N}_{>0}$，可知$\sum_{k=m+1}^n a_k \cdot \binom{k}{k-m+1}h^{k-m+1}$ 是有理数.\par
    而$a_m$是无理数，所以$mha_m$是无理数.故$q_h(j)$的$j^{m-1}$次项系数$c_{m-1}$是无理数.\par
    $q_h(j)$是次数为$n-1$的多项式且$j^{m-1}$的系数为无理数.根据假设，对于$\forall\, h>0$，$q_h(j) \mod 1$在$[0,1)$中均匀分布，即$\frac{1}{J}\sum_{j=0}^{J-1}{\overline{b_{j+h}}b_j}\to 0\;(\forall\, h\in\mb{N}_{>0})$. \par
    因此，$\frac{1}{J}\sum_{j=0}^{J-1}{\eu^{2k\pi \iu p(j)}}\to 0\;(\forall\, k\in\mb{Z}\setminus\{0\})$，$p(x)\mod 1$在$[0,1)$中均匀分布.\par 
    由归纳法可知上述命题对$n\in\mb{N}_{n>0}$均成立.
\end{proof}


\section{回复定理}

下面学习一些回复定理并使用回复定理回答一些问题.\par
\begin{theorem}{Poincare回复定理}
    设$(X,\mcal{B},m)$为概率测度空间，$T:X\to X$是保测映射，$A\in\mcal{B},m(A)>0$，那么存在$F\subseteq A,m(A\setminus F)=0$，使得$\forall\, x\in F$，从$x$出发的轨迹$\{x,Tx,T^2x,\cdots,T^nx,\cdots\}$无限多次回到$A$中(更进一步，回到$F$中）.
\end{theorem}
\begin{proof}
    记$B_1=T^{-1}A\cup T^{-2}A\cup\cdots\cup T^{-n}A\cup\cdots$，即$\forall\, x\in B,\exists\, n\in\mb{N},T^nx\in A$.\par
    $B_2=T^{-2}A\cup T^{-3}A\cup\cdots\cup T^{-(n+1)}A\cup\cdots$，即$\forall\, x\in B_2,\exists\, n\in\mb{N},n\geq 2,T^nx\in A$.\par
    $\quad\vdots$\par
    $B_k=T^{-k}A\cup T^{-(k+1)}A\cup\cdots$，以此类推.\par
    令$B_0=A\cup T^{-1}A\cup\cdots\cup T^{-n}A\cup\cdots$.\par
    则$B_0\supseteq B_1\supseteq B_2\supseteq\cdots\supseteq B_k\supseteq\cdots$，且$T^{-1}B_k=B_{k+1},m(B_k)=m(B_{k+1})$.\par
    令$B_\infty=\bigcap_{i=0}^\infty{B_i}=\bigcap_{i=1}^\infty{B_i}$，那么$m(B_\infty)=m(B_k)$.\par
    则$\forall\, x\in B_\infty$，$x$出发的轨道无限次进入集合$A$.\par
    令$F=B_\infty\cap A$，由$m(B_\infty)=m(B_0),B_0=A\cup T^{-1}A\cup T^{-2}A\cup\cdots\supseteq A$，得$m(F)=m(A)$.\par
    下面证明$\forall\, x\in F$，$x$出发的轨道无限次进入集合$F$.\par
    $\forall\, x\in F$，存在子列$n_1<n_2<\cdots<n_k<\cdots$，使得$T^{n_1}x\in A,T^{n_2}x\in A,\cdots,T^{n_k}x\in A,\cdots$.\par
    对于$T^{n_1}x$，$T^{n_i}x=T^{n_i-n_1}\left(T^{n_1}x\right)\in A\;(\forall\, i\geq 2)$，因此$T^{n_1}x\in F$.\par
    类似地，$T^{n_2}x,T^{n_3}x,\cdots\in F$.
\end{proof}

下面的回复定理跳出了概率测度空间的范畴，它定义在度量空间上.
\begin{theorem}{Birkhoff回复定理}
    设$(X,d)$为紧致度量空间，$T:X\to X$是连续映射，存在$x\in X$，以及自然数子列$\{n_i\}_{i\in\mb{N}}$，使得$T^{n_i}x\to x(i\to\infty)$
\end{theorem}

\begin{definition}{$\omega-$极限集和$\alpha-$极限集}
    $\forall\, x\in X$，记$\omega(x)$为序列$Tx,T^2x,\cdots,T^nx,\cdots$所有极限点构成的集合，称为$x$的\textbf{$\bm{\omega}$--极限集}.\par
    如果$T$为同胚，则记$\alpha(x)$为序列$T^{-1}x,T^{-2}x,\cdots,T^{-n}x,\cdots$所有极限点构成的集合，称为$x$的\textbf{$\bm{\alpha}$--极限集}.
\end{definition}

\begin{proposition}{$\omega$--极限集的性质}
    \begin{enumerate}
        \item $\forall\, x\in X,\omega(x)\neq\varnothing$.
        \item $\omega(x)$是闭集(紧集).
        \item $\omega(x)$是$T-$不变的，即$T\omega(x)=\omega(x)$.
    \end{enumerate}
\end{proposition}
\begin{proof}
    $\forall\, y\in T\omega(x)$，$\exists\, z\in\omega(x)$，使得$y=Tz$.\par
    设$z=\lim_{i\to\infty}T^{n_i}x\;(n_1<n_2<\cdots<n_k<\cdots)\Rightarrow y=\lim_{i\to\infty}T^{n_i+1}x\Rightarrow y\in\omega(x)$.\par
    因此$T\omega(x)\subseteq \omega(x)$.\par
    $\forall\, y\in\omega(x)$，$\exists\, n_1<n_2<\cdots<n_k<\cdots$，使得$\lim_{i\to\infty}{T^{n_i}x}=y$，则$\lim_{i\to\infty}T \left(T^{n_i-1}(x)\right)=y$.\par
    由$X$的紧致性，取$\left\{T^{n_i-1}(x)\right\}$的收敛子列$\left\{T^{m_i}(x)\right\}$，并记$T^{m_i}(x)\to z$.\par
    那么由连续性$y=Tz\Rightarrow y\in T\omega(x)$.\par
    因此$\omega(x)\subseteq T\omega(x)$.
\end{proof}

\begin{definition}{极小系统}
    设$(X,d)$是紧致度量空间，$T:X\to X$是连续映射，如果$\forall\, x\in X,\left\{x,Tx,\cdots,T^nx,\cdots\right\}$在$X$中稠密，那么称$T:X\to X$为\textbf{极小系统}.\par
    特别地，若$T:X\to X$是同胚，如果$\forall\, x\in X, \left\{\cdots,T^{-2}x,T^{-1}x,x,Tx,T^2x,\cdots\right\}$在$X$中稠密，那么称$T:X\to X$为\textbf{极小系统}.
\end{definition}

\begin{lemma}{Zorn引理}
    设$\mcal{F}$是一个集合，$\prec $是$\mcal{F}$上的一个偏序关系.如果对于任意一个全序子集$A\subseteq\mcal{F}$，可以找到元素$x$使得$\forall\, y\in A,x<y$，那么$\mcal{F}$中存在极小元.
\end{lemma}

\begin{theorem}
    设$(X,d)$为紧致度量空间，$T:X\to X$是连续映射，那么存在$X$中的非空闭集$A$，使得$TA\subseteq A$，且$T:A\to A$是极小系统.
\end{theorem}














